\section{Introduction}\label{sec:introduction}
Industrial anomaly detection (AD) is typically deployed under an asymmetric data regime: normal examples are available, while abnormal examples are rare, diverse, and costly to annotate. This makes few-shot AD practically important. In this setting, an inspection system receives only $k$ normal support images for a target category, often with $k\in\{1,2,4,8\}$, and must detect both image-level and pixel-level defects at test time.

Recent methods have made this regime much more effective by using frozen visual foundation models. DINOv2~\cite{dinov22023} patch features can be stored in a normal memory bank and queried by nearest-neighbor distances, as in PatchCore-style methods and AnomalyDINO~\cite{patchcore2021,anomalydino2024}. Alternatively, frozen patch features can be compressed into a normal PCA/subspace model, as in SubspaceAD~\cite{subspacead2026}. These approaches provide strong anomaly ranking without training a large task-specific backbone.

However, ranking is not the same as reliable decision support. A raw nearest-neighbor distance or PCA residual can produce high AUROC while still being poorly calibrated: a score threshold may work on one class, but the same numerical score may not carry the same false-alarm meaning under another class, dataset, or corruption. This matters in industrial deployment because operators need to know not only which samples are suspicious, but also whether an alarm probability is trustworthy. Prior work has already shown that DINOv2-based few-shot AD can be poorly calibrated and adversarially fragile~\cite{khan2025calibration}. Reliability is therefore not a cosmetic post-processing issue; it is part of the deployment problem.

We address this gap by separating anomaly ranking from reliability estimation. The ranking path remains a frozen DINOv2 PCA/subspace detector: patch residuals are aggregated into a raw anomaly score used for AUROC, AP, and localization. A separate reliability path maps this evidence into calibrated or probability-like outputs. We call the resulting framework \emph{Conformal Reliability Routing} (CRR). The main CRR route uses leave-one-image-out (LOIO) conformal calibration over the few normal support images. Each support image is scored as if held out, producing normal nonconformity values; a test image then receives a conformal p-value and a reliability score $1-p$. This view is not a supervised anomaly posterior, but it is statistically interpretable and derived without target anomaly labels.

Each ingredient we build on---DINOv2 features, PCA residuals, Platt-family calibration, and conformal anomaly detection---has an established literature, and our contribution begins after them. What that literature leaves open is a structural question about the few-shot regime itself: a conformal alarm calibrated on $k$ support images cannot fire below $\alpha=1/(k{+}1)$ at all, and under corruption shift its false-alarm promise can silently invert. The central contribution of this paper is that pair of problem and mechanism---making the resolution floor and its failure directions exact and checkable, and then introducing a source-validated cross-category scheme (SC3R) that alarms below the floor with controlled false alarms and no target anomaly labels. The decoupled detector and the calibration benchmark around them are the supporting substrate: a protocol under which these claims can be tested without conflating ranking and reliability.

The paper makes four contributions, the first two central:
\begin{itemize}
    \item We make the few-shot resolution limit of conformal alarms explicit through an attainable-alpha analysis: with $k$ supports no alarm can fire below $\alpha=1/(k{+}1)$, and empirical false-alarm rates must be read against the attainable grid, not the nominal level. Exact discrete uniformity tests of normal p-values then localize when and in which direction the false-alarm promise fails under corruption shift---conservative on VisA, anti-conservative on corrupted MVTec.
    \item We introduce SC3R, a source-validated cross-category thresholding scheme that pools normal images from other categories to unlock operating points below the target-only attainable-alpha floor. Under a pre-registered decision gate on all 15 MVTec classes, matched-condition SC3R tracks nominal false-alarm rates almost exactly with substantial power where the target-only anchor is structurally silent, and hierarchical class--seed bootstrap intervals for its power gain exclude zero for every corruption condition. The result replicates on all 12 VisA classes, and cross-dataset source archives (MVTec sources for VisA targets) transfer conservatively: reduced power with false alarms well below nominal.
    \item Supporting these claims, we formulate a decoupled few-shot AD pipeline in which frozen DINOv2 PCA/subspace residuals provide anomaly ranking while a separate label-free reliability layer---vector/shift-aware calibration and LOIO conformal p-values~\cite{hennhofer2024loio} computed only from the $k$ normal supports---provides decision signals, with full-benchmark corruption evidence on all 12 VisA and all 15 MVTec classes that the conformal route improves calibration over Platt-family calibrators and over standard scalar calibrators (temperature scaling, isotonic regression, histogram binning) under paired Wilcoxon tests while preserving the raw ranking.
    \item We report accuracy-storage, transfer calibration, pixel-level, robustness, prevalence-stress, and routing diagnostics, including negative findings that prevent overclaiming: CRR is not MVTec AUROC SOTA, is not adversarially robust, and ECE on conformal-derived scores is strongly prevalence-sensitive.
\end{itemize}
